\documentclass[11pt]{article}
\usepackage[T1]{fontenc}
\usepackage[utf8]{inputenc}
\usepackage{mathptmx}
\usepackage[margin=1in]{geometry}
\usepackage[hidelinks]{hyperref}
\usepackage{enumitem}
\usepackage{microtype}
\usepackage{tabularx}
\usepackage[dvipsnames]{xcolor}
\definecolor{ReviewerGreen}{HTML}{008000}
\newcommand{\rvcomment}[1]{{\color{ReviewerGreen}#1}}
\newcommand{\concern}[1]{\textbf{\color{ReviewerGreen}#1}}
\setlength{\parindent}{0pt}
\setlength{\parskip}{6pt}
\setlist[itemize,1]{leftmargin=1.6em,labelsep=0.6em,itemsep=3pt,topsep=3pt,parsep=0pt,
                    label=\raisebox{-0.15ex}{\large\textbullet}}
\setlist[itemize,2]{leftmargin=1.4em,labelsep=0.5em,itemsep=2pt,topsep=3pt,parsep=0pt,
                    label=\textendash}
\pagestyle{plain}
\begin{document}
\noindent\begin{tabularx}{\textwidth}{@{}l@{\hspace{1.1em}}X@{}}
\textbf{Original Manuscript ID:} & Access-2026-38801 \\[3pt]
\textbf{Original Article Title:} & ``The Price of De-correlation: Quantifying the Agreement-Cost Trade-off in Heterogeneous Large Language Model Ensembles''
\end{tabularx}

\vspace{18pt}

\noindent\begin{tabularx}{\textwidth}{@{}l@{\hspace{1.1em}}X@{}}
\textbf{To:} & IEEE Access Editor \\[3pt]
\textbf{Re:} & Response to reviewers
\end{tabularx}

\vspace{16pt}
Dear Editor,

Thank you for allowing a resubmission of our manuscript, with an opportunity to address the reviewers' comments.

We are uploading (a) our point-by-point response to the comments (below) (response to reviewers, under ``Author's Response Files''), (b) an updated manuscript with yellow highlighting indicating changes (as ``Highlighted PDF''), and (c) a clean updated manuscript without highlights (``Main Manuscript'').

The revision adds two substantial pieces of evidence. First, we now measure the correlation between agents' errors directly rather than inferring it from label agreement, which is what the term ``de-correlation'' in our title asserts. Second, we ran a pre-registered replication of the entire protocol on a substantially different domain, 537 expert clinical-decision items, and report it whatever the outcome. Every number below is regenerated by the released code from the frozen data.

\vspace{14pt}
Best regards, \\[22pt]
Samir Chincholikar \\
Robin Chawla

\newpage

\vspace{12pt}
\concern{Reviewer \#1, Concern \#1}

\concern{Reviewer comment:}

\rvcomment{Reduced agreement is not shown to improve accuracy or any decision-relevant outcome. Link the agreement-cost frontier to calibration, abstention, robustness, or tail-error reduction.}

\textbf{Author response:}

\begin{itemize}

\item We agree that reduced agreement alone does not establish a downstream benefit.

\item We added a selective-prediction analysis using mean agent conviction to rank ensemble decisions. AURC is 0.553 for HOM, 0.568 for HET-LITE, and 0.554 for HET. The HOM-HET difference is -0.001 with a 95\% confidence interval of [-0.059, 0.060]. The corresponding two- way interval is [-0.113, 0.115]. Both intervals contain zero.

\item Restricting evaluation to the most confident half of decisions does not materially reduce selective risk in any configuration. We therefore claim no improvement in calibration or abstention quality.

\item The exploratory consensus analysis separately evaluates unanimous incorrect decisions and is addressed under Reviewer \#1, Concern \#4.

\end{itemize}

\textbf{Author action:}

\begin{itemize}

\item Section V-I, ``Cross-Domain Replication: Clinical Knowledge,'' subsection ``Selective Prediction,'' was added to report the risk-coverage analysis and clustered AURC contrast.

\item Section V-D, ``The Agreement-Cost Frontier,'' was revised to state explicitly that the frontier describes the cost of greater behavioral diversity and does not establish improved downstream decision quality.

\end{itemize}

\vspace{12pt}
\concern{Reviewer \#1, Concern \#2}

\concern{Reviewer comment:}

\rvcomment{The tier-matched control covers only two providers and no xAI equivalent. A fully tier-matched multi-provider design or model-level random effects is needed.}

\textbf{Author response:}

\begin{itemize}

\item A fully three-provider tier-matched design remains unavailable because an equivalent efficient-tier xAI model was not available.

\item We therefore added an additional capability comparison in the clinical replication. Within the fully heterogeneous configuration, one same-provider pair and one cross-provider pair have nearly identical mean member accuracy, 0.951 and 0.955 respectively.

\item Their error correlations are 0.954 and 0.607, a difference of 0.347 with a 95\% confidence interval of [0.207, 0.507].

\item We describe this as approximate capability matching. Provider remains confounded with other model-specific characteristics and is not identified as an independent causal factor.

\end{itemize}

\textbf{Author action:}

\begin{itemize}

\item Section V-I, ``Cross-Domain Replication: Clinical Knowledge,'' was revised to add the approximate capability-matched provider comparison.

\item Section V-J, ``Robustness: Capability-Matched Control,'' retains the pre-registered financial tier-matched analysis and states explicitly that matching is approximate and covers only two providers.

\item Section VII, ``Threats to Validity,'' subsection ``Heterogeneity vs. capability confound,'' was revised to retain the unresolved provider and capability limitations.

\end{itemize}

\vspace{12pt}
\concern{Reviewer \#1, Concern \#3}

\concern{Reviewer comment:}

\rvcomment{Findings come from one financial task, one prompt, three labels, and fixed five-agent ensembles. Test other domains, ensemble sizes, aggregation rules, and label structures.}

\textbf{Author response:}

\begin{itemize}

\item We added a pre-registered replication on a substantially different clinical-knowledge task using 537 questions from the MMLU professional-medicine and clinical-knowledge subsets.

\item The response structure changes from three financial labels to four answer options with one correct answer. The five-agent structure, three configurations, seeding procedure, temperature, number of runs, and estimators remain unchanged.

\item The directional ordering reproduces. Agreement decreases by delta-kappa = 0.072 with a 95\% confidence interval of [0.040, 0.104], compared with 0.336 in finance. Error correlation decreases by delta-phi = 0.393 with a 95\% confidence interval of [0.325, 0.459], compared with 0.454 in finance.

\item We have not varied ensemble size or aggregation rule and now identify those remaining dimensions explicitly as limitations.

\end{itemize}

\textbf{Author action:}

\begin{itemize}

\item Section V-I, ``Cross-Domain Replication: Clinical Knowledge,'' was added to report the pre- registered clinical replication.

\item Section VII, ``Threats to Validity,'' subsection ``Scope and generalizability,'' was revised to state which dimensions were replicated and which remain fixed.

\item The replication pre-registration and freeze receipt were added to the released reproducibility materials.

\end{itemize}

\vspace{12pt}
\concern{Reviewer \#1, Concern \#4}

\concern{Reviewer comment:}

\rvcomment{The 1.84x unanimous-incorrect finding and HET-LITE recommendation are post hoc, lack CIs/p-values, and ignore multiplicity. Treat as hypothesis-generating only.}

\textbf{Author response:}

\begin{itemize}

\item We now treat this analysis explicitly as exploratory.

\item The HOM-to-HET unanimous-incorrect ratio is 1.84. The equity-clustered 95\% confidence interval is [1.42, 2.52]. Under the two-way equity and date resampling scheme, the interval is [1.13, 3.06]. Both exclude unity.

\item The comparison and statistic were selected after inspection of the data, multiple related comparisons were available, and no multiplicity adjustment is applied. The manuscript therefore does not treat these intervals as confirmatory evidence.

\item HET-LITE is no longer presented as a recommended design. Its observed cost efficiency is treated only as a hypothesis for prospective evaluation.

\end{itemize}

\textbf{Author action:}

\begin{itemize}

\item Section V-H, ``Exploratory Analysis of Consensus Errors,'' was revised to add equity-clustered and two-way uncertainty for the unanimous-incorrect ratio.

\item Section V-H was also revised to add the post hoc and multiplicity caveats and to remove the HET-LITE deployment recommendation.

\item Section VI, ``Discussion,'' was revised so that the HET-LITE result is described as hypothesis- generating rather than deployment guidance.

\end{itemize}

\vspace{12pt}
\concern{Reviewer \#1, Concern \#5}

\concern{Reviewer comment:}

\rvcomment{Seed support is best-effort, and the diagnostic is strongest only for HOM. Heterogeneous cross-run checks are confounded by repeated same-model draws. Add verification or sensitivity analysis.}

\textbf{Author response:}

\begin{itemize}

\item We added a sensitivity analysis that separates pairwise agreement according to whether two agents use the same model.

\item Same-model agreement is 0.800 in HOM, 0.794 in HET-LITE, and 0.862 in HET. Different-model agreement is 0.455 in HET-LITE and 0.490 in HET.

\item Same-model agreement remains high across configurations, while different-model pairs exhibit substantially lower agreement.

\item We also clarify that distinct per-agent seeds remove the demonstrated shared-seed artifact but do not establish statistical independence of model responses.

\end{itemize}

\textbf{Author action:}

\begin{itemize}

\item Section V-E, ``Shared-Seed Dependence and Measurement Validation,'' subsection ``Sensitivity: agreement by pair type,'' was added.

\item Section V-E, subsection ``What unique seeds do and do not establish,'' was added to state explicitly that distinct seeds do not prove independence.

\item Figure 2 was revised to describe the change as assignment of a distinct reproducible seed to each agent rather than restoration of statistical independence.

\end{itemize}

\vspace{12pt}
\concern{Reviewer \#1, Concern \#6}

\concern{Reviewer comment:}

\rvcomment{Table 4 and the text disagree on costs. Clarify per-decision cost, token accounting, and pricing assumptions.}

\textbf{Author response:}

\begin{itemize}

\item We clarified the cost calculation. For each call, metered prompt and completion tokens are priced using the rates in the model registry. Costs are summed across the five agents in an ensemble decision and averaged across cells.

\item The resulting costs are \$6.48 x 10\textasciicircum{}-4 for HOM, \$1.52 x 10\textasciicircum{}-3 for HET-LITE, and \$2.84 x 10\textasciicircum{}-3 for HET. The HOM-to-HET cost ratio is approximately 4.4x.

\item Prices are fixed at those applicable during data collection. No discount, batching, or caching adjustment is applied. Failed requests returning no billable completion are excluded.

\end{itemize}

\textbf{Author action:}

\begin{itemize}

\item Section V-D, ``The Agreement-Cost Frontier,'' was revised to state the token accounting, aggregation, and pricing assumptions explicitly.

\item Table 4 reports the corresponding per-ensemble-decision costs for all three configurations.

\end{itemize}

\vspace{12pt}
\concern{Reviewer \#2, Concern \#1}

\concern{Reviewer comment:}

\rvcomment{Simplify index terms to 3 - 5 words only.}

\textbf{Author response:}

\begin{itemize}

\item Done. The list has been reduced to five index terms.

\end{itemize}

\textbf{Author action:}

\begin{itemize}

\item The manuscript Index Terms were revised to: algorithmic monoculture, inter-rater agreement, large language models, model heterogeneity, and multi-agent systems. They are given in alphabetical order, as the IEEE Editorial Style Manual requires.

\end{itemize}

\vspace{12pt}
\concern{Reviewer \#2, Concern \#2}

\concern{Reviewer comment:}

\rvcomment{No clear explanation about fair comparison with current state-of-the-art approach.}

\textbf{Author response:}

\begin{itemize}

\item This article does not propose a new ensemble algorithm, so a direct task-performance competition with existing ensemble methods is not the primary comparison.

\item Existing methods such as mixture-of-agents, multi-agent debate, and diversity-based ensemble selection are generally evaluated through downstream performance. The present study instead measures how ensemble composition changes output agreement, error dependence, and inference cost under fixed task conditions.

\item The contribution is therefore a measurement framework that separates output agreement from error correlation, controls the shared-seed artifact, quantifies uncertainty, and incorporates inference cost.

\end{itemize}

\textbf{Author action:}

\begin{itemize}

\item Section V-L, ``Relation to Existing Ensemble Methods,'' was added to explain the appropriate basis of comparison with existing ensemble approaches.

\item Section II-A, ``LLM Ensembles and Mixture-of-Agents,'' provides the corresponding literature context for diversity-based and heterogeneous ensemble approaches.

\end{itemize}

\vspace{12pt}
\concern{Reviewer \#3, Concern \#1}

\concern{Reviewer comment:}

\rvcomment{The present experiment confounds provider, model identity, capability, training data, architecture, and alignment. Add same-provider/different-model controls, more than one model per provider, balanced provider mixtures, and preferably a factorial analysis; otherwise rewrite all causal/provider-specific claims as configuration-specific observational findings.}

\textbf{Author response:}

\begin{itemize}

\item We agree that provider is not experimentally separable from model identity, training data, architecture, alignment procedure, and other model-specific characteristics in the primary design.

\item We therefore interpret the provider comparisons as configuration-specific observational findings rather than isolated causal provider effects.

\item The pre-registered capability-matched financial control preserves approximately 94\% of the primary agreement contrast. The clinical replication additionally provides an approximate capability-matched provider comparison with nearly identical member accuracy.

\item These analyses narrow the capability confound but do not identify provider as an independent causal factor.

\end{itemize}

\textbf{Author action:}

\begin{itemize}

\item Section V-F, ``Provider-Pair Agreement,'' was rewritten as a configuration-specific observational analysis and now states explicitly that provider is bundled with other model characteristics.

\item Section V-I, ``Cross-Domain Replication: Clinical Knowledge,'' was revised to include the approximate capability-matched same-provider versus cross-provider comparison.

\item Section V-J, ``Robustness: Capability-Matched Control,'' retains the pre-registered financial capability control.

\item Section VII, ``Threats to Validity,'' subsection ``Heterogeneity vs. capability confound,'' was revised to state that provider cannot be isolated causally in the present design.

\end{itemize}

\vspace{12pt}
\concern{Reviewer \#3, Concern \#2}

\concern{Reviewer comment:}

\rvcomment{Recompute uncertainty while accounting for both repeated observations within equities and shared analysis dates/common market conditions. Justify the bootstrap unit formally, evaluate two-way equity/date clustering or an appropriate multiway resampling strategy, and regenerate all main and exploratory confidence intervals accordingly.}

\textbf{Author response:}

\begin{itemize}

\item Equity remains the pre-registered primary resampling unit because repeated dates are nested within securities. Analysis date is treated as a second crossed source of dependence.

\item For the primary contrast, the point estimate is 0.3363 under all resampling schemes. The 95\% intervals are [0.3035, 0.3689] under equity clustering, [0.2989, 0.3656] under date clustering, and [0.2786, 0.3877] under two-way resampling. All exclude zero.

\item Two-way sensitivity analysis is also reported for the inferential exploratory contrasts. The unanimous-incorrect ratio has an interval of [1.13, 3.06], the pile-on contrast has [0.140, 0.232], and the AURC contrast has [-0.113, 0.115].

\item The error-correlation contrasts also retain their conclusions under two-way resampling, as reported under Reviewer \#3, Concern \#3.

\end{itemize}

\textbf{Author action:}

\begin{itemize}

\item Section III-D, ``Uncertainty Quantification,'' was revised to justify equity clustering formally and to define the date-clustered and two-way resampling procedures.

\item Table 5 was added to report the primary HOM-HET contrast under equity, date, and two-way resampling.

\item Section V-C, ``Agreement, Error Correlation, and Decorrelation,'' was revised to report two-way sensitivity intervals for both error-correlation specifications.

\item Section V-H, ``Exploratory Analysis of Consensus Errors,'' was revised to report two-way intervals for the unanimous-incorrect ratio and pile-on contrast.

\item Section V-I, subsection ``Selective Prediction,'' was revised to report the two-way AURC interval.

\end{itemize}

\vspace{12pt}
\concern{Reviewer \#3, Concern \#3}

\concern{Reviewer comment:}

\rvcomment{Fleiss' kappa measures agreement, not correlation. Either remove ``de-correlation''/``less-correlated judgments'' throughout the title, abstract, figures, discussion, and conclusion, or add actual dependence/error-correlation analyses that justify those terms. Explicitly distinguish output agreement, error correlation, sampling independence, model diversity, and provider diversity.}

\textbf{Author response:}

\begin{itemize}

\item We took the second option and now measure error dependence directly.

\item For every pair of agents within a cell, we compute phi between binary error indicators. Error correlation decreases from 0.984 for HOM to 0.742 for HET-LITE and 0.530 for HET. The HOM-HET contrast is 0.454 with a 95\% confidence interval of [0.376, 0.534]. The two-way interval is [0.328, 0.586].

\item We repeat the analysis with HOLD counted as an error so that all configurations use the same evaluated subset. The corresponding contrast is 0.309 with a 95\% confidence interval of [0.262, 0.354]. The two-way interval is [0.227, 0.387].

\item The manuscript now distinguishes output agreement, error correlation, sampling independence, model diversity, and provider diversity. Decorrelation is defined specifically as a reduction in error correlation measured by phi, while kappa is retained as the agreement measure.

\end{itemize}

\textbf{Author action:}

\begin{itemize}

\item Section V-C, ``Agreement, Error Correlation, and Decorrelation,'' was added to define the five concepts and report the phi analysis under both treatments of HOLD.

\item Section V-A was retitled ``Heterogeneity Reduces Ensemble Agreement (Primary)'' so that the kappa result is no longer described as decorrelation.

\item The Abstract was revised to report the phi contrast alongside the kappa contrast.

\item Section VI, ``Discussion,'' and Section VIII, ``Conclusion,'' were revised so that claims about decorrelation refer to error dependence rather than categorical agreement.

\end{itemize}

\vspace{12pt}
\concern{Reviewer \#3, Concern \#4}

\concern{Reviewer comment:}

\rvcomment{HOLD-induced abstention creates different evaluation subsets across configurations, so the reported hit rates are not directly comparable. Perform a proper selective-prediction/coverage analysis, use paired and clustered inference, provide uncertainty for unanimous-error and pile-on statistics, and stop using overlapping marginal Wilson intervals as the main inferential evidence.}

\textbf{Author response:}

\begin{itemize}

\item The selective-prediction analysis reports a HOM-HET AURC difference of -0.001 with a 95\% equity-clustered confidence interval of [-0.059, 0.060] and a two-way interval of [-0.113, 0.115].

\item The unanimous-incorrect ratio is 1.84 with an equity-clustered interval of [1.42, 2.52] and a two-way interval of [1.13, 3.06].

\item For pile-on, the HOM-HET difference is 0.189 with an equity-clustered interval of [0.160, 0.216] and a two-way interval of [0.140, 0.232].

\item The error-correlation analysis is also reported with HOLD counted as an error, which fixes the evaluated subset across configurations.

\item Wilson intervals are retained only for the descriptive secondary accuracy analysis and are no longer the primary inferential evidence.

\end{itemize}

\textbf{Author action:}

\begin{itemize}

\item Section V-I, subsection ``Selective Prediction,'' was added to report risk-coverage analysis and paired clustered AURC inference.

\item Section V-H, ``Exploratory Analysis of Consensus Errors,'' was revised to report clustered and two-way uncertainty for unanimous-incorrect decisions and pile-on.

\item Section V-C, ``Agreement, Error Correlation, and Decorrelation,'' was revised to report the common-subset error-correlation analysis with HOLD counted as an error.

\item Section V-G, ``Secondary Directional-Accuracy Analysis,'' now treats the accuracy results as secondary and does not rely on marginal Wilson intervals as the principal inferential evidence.

\end{itemize}

\vspace{12pt}
\concern{Reviewer \#3, Concern \#5}

\concern{Reviewer comment:}

\rvcomment{Clarify that unique seeds do not prove statistical independence; provide complete API parameters and versions, exact prompts/system messages, seed-support evidence, retry records, total attempted versus successful API calls, exact preregistration timestamp evidence, and archived immutable code/data. Correct Figure 3 so that the shared-seed pilot, independent-seed pilot, and full-study quantities compare the same estimand.}

\textbf{Author response:}

\begin{itemize}

\item The manuscript now states explicitly that distinct per-agent seeds remove the demonstrated shared-seed artifact but do not prove statistical independence.

\item The confirmatory grid contains 54,000 valid planned responses. Completing it required 94,566 API requests. Of the 40,566 unsuccessful requests, 40,544 were rate-limit or quota responses and 22 were transport or schema failures.

\item Every request attempt is retained in the released log. The reproducibility materials contain the request settings, prompts, prompt hashes, model identifiers, prices, experimental configuration, pre-registration records, and archived model responses.

\item Figure 3 now contains only the two like-for-like pilot HOM-HET agreement contrasts. The contrast is 0.485 under a shared seed and 0.113 under distinct per-agent seeds.

\item The full-study within-minus-cross-run diagnostic is reported separately in Table 6 and is no longer mixed with the Figure 3 estimand.

\end{itemize}

\textbf{Author action:}

\begin{itemize}

\item Section III-C, ``Independent-Draw Estimand,'' was revised to state the seeding construction and the limits of seed-based reproducibility.

\item Section IV-C, ``Models,'' records the provider and model identifiers and fixed inference settings.

\item Section IV-D, ``Pre-Registration, Freezing, and Reproducibility,'' documents the frozen design, hashes, archived dataset, and deterministic analysis pipeline.

\item Section IV-E, ``Collection Provenance and Retry Accounting,'' was added to report the 94,566 attempted requests, 54,000 valid outputs, failure breakdown, and retained attempt logs.

\item Section V-E, ``Shared-Seed Dependence and Measurement Validation,'' was revised to distinguish the shared-seed artifact from statistical independence.

\item Figure 2 was revised to describe the distinct per-agent seed construction without claiming independence.

\item Figure 3 was reduced to the two directly comparable pilot HOM-HET agreement contrasts.

\item Table 6 reports the full-study within versus cross-run residual dependence diagnostic separately.

\end{itemize}

\vspace{12pt}
\concern{Reviewer \#3, Concern \#6}

\concern{Reviewer comment:}

\rvcomment{Replicate the main experiment on at least one substantially different domain/task and preferably multiple prompt formats, ensemble sizes, and model mixtures. The present single BUY/HOLD/SELL financial task-whose accuracy is itself near chance and whose label marginals vary strongly by configuration-is insufficient to establish a general principle about heterogeneous LLM ensembles.}

\textbf{Author response:}

\begin{itemize}

\item We added a pre-registered replication using 537 MMLU professional-medicine and clinical- knowledge questions. Chance performance is 25\%, and the evaluated models perform substantially above that level.

\item The response structure contains four options rather than the three BUY, HOLD, and SELL labels. Choice order is shuffled deterministically to remove positional imbalance in the published answer key.

\item The main ordering reproduces. Agreement decreases by delta-kappa = 0.072 with a 95\% confidence interval of [0.040, 0.104], while error correlation decreases by delta-phi = 0.393 with a 95\% confidence interval of [0.325, 0.459].

\item Agreement changes substantially more across the two domains than error correlation. We therefore do not claim that the magnitude is invariant across domains.

\item Prompt format, five-agent ensemble size, majority-vote aggregation, and the tested model mixtures remain fixed and are identified explicitly as limitations.

\end{itemize}

\textbf{Author action:}

\begin{itemize}

\item Section V-I, ``Cross-Domain Replication: Clinical Knowledge,'' was added to report the pre- registered 537-item clinical replication.

\item Section VII, ``Threats to Validity,'' subsection ``Scope and generalizability,'' was revised to state explicitly that the result has been evaluated across two substantially different tasks while prompt format, ensemble size, aggregation rule, and model mixtures remain fixed.

\item The clinical replication pre-registration and freeze receipt were added to the released reproducibility materials.

\end{itemize}

\end{document}
